% Source: Weinberg GR, physical PDF pp. 187--189
%         (printed pp. 163--165).
% Coverage: Section 7.5, The Cauchy Problem; equations (7.5.1)--(7.5.3).
% Figures/tables/footnotes: none.
% Status: source-reviewed and compile-clean.
% Uncertainties: none.
% MODERNIZATION: Einstein-equation and constraint-propagation signs follow
% the target G_{\mu\nu}=+8\pi G T_{\mu\nu} convention.

\subsection{The Cauchy Problem}\label{sec:7.5}

We can gain further insight into the mathematical content of Einstein's
equations by applying them to the traditional initial-value problem of
Cauchy. Suppose that we are given \(g_{\mu\nu}\) and
\(\partial_0g_{\mu\nu}\) everywhere on the ``plane'' \(x^0=t\). If we could
extract from the field equations a formula for
\(\partial_0^2g_{\mu\nu}\) everywhere at \(x^0=t\), we could then compute
\(g_{\mu\nu}\) and \(\partial_0g_{\mu\nu}\) at a time \(x^0=t+\delta t\),
and by continuing this process \(g_{\mu\nu}\) could be computed for all
\(x^i\) and \(x^0\).

At first sight this looks feasible, because we need 10 second derivatives,
and there are 10 field equations. But let us look more closely at the
left-hand side
\[
  G^{\mu\nu}\equiv R^{\mu\nu}-\frac12g^{\mu\nu}R
\]
of the field equations. The Bianchi identities \eqref{eq:6.8.4} tell us that
\[
  \partial_0G^{\mu0}
  \equiv
  -\partial_iG^{\mu i}
  -\Gamma^\mu{}_{\nu\lambda}G^{\lambda\nu}
  -\Gamma^\nu{}_{\nu\lambda}G^{\mu\lambda}.
\]
The right-hand side contains no time derivatives higher than
\(\partial_0^2\), so neither does the left-hand side, and therefore
\(G^{\mu0}\) contains no time derivatives higher than \(\partial_0\). Thus
we cannot learn anything about the time evolution of the gravitational field
from the four equations
\begin{equation}
  G^{\mu0}=8\pi G T^{\mu0}.
  \tag{7.5.1}\label{eq:7.5.1}
\end{equation}
Rather, these equations must be imposed as constraints on the initial data,
that is, on \(g_{\mu\nu}\) and \(\partial_0g_{\mu\nu}\) at \(x^0=t\).

This leaves as ``dynamical'' equations only the other six Einstein equations
\begin{equation}
  G^{ij}=8\pi G T^{ij}.
  \tag{7.5.2}\label{eq:7.5.2}
\end{equation}
When we solve these equations for the 10 second derivatives
\(\partial_0^2g_{\mu\nu}\), we must encounter a fourfold ambiguity, which of
course we could not have hoped to escape, since it is always possible to make
coordinate transformations that leave \(g_{\mu\nu}\) and
\(\partial_0g_{\mu\nu}\) unchanged at \(x^0=t\) but that do alter
\(g_{\mu\nu}\) everywhere else. To be more specific, what we find is that
Eq.~\eqref{eq:7.5.2} determines the six \(\partial_0^2g^{ij}\), but leaves
the other four derivatives \(\partial_0^2g^{\mu0}\) indeterminate. This
ambiguity can be removed by imposing four coordinate conditions that fix the
coordinate system. For instance, if we adopt the harmonic coordinate
condition discussed in the last section, the second time derivative of
\(\sqrt{-g}\,g^{\mu0}\) can be determined by differentiating
Eq.~\eqref{eq:7.4.6} with respect to time:
\begin{equation}
  \partial_0^2\left(\sqrt{-g}\,g^{\mu0}\right)
  =
  -\partial_0\partial_i\left(\sqrt{-g}\,g^{\mu i}\right),
  \tag{7.5.3}\label{eq:7.5.3}
\end{equation}
and the 10 equations \eqref{eq:7.5.2} and \eqref{eq:7.5.3} suffice to
determine the second time derivatives of all \(g_{\mu\nu}\).

When the initial-value problem is solved in this way, the constraints
\eqref{eq:7.5.1} on the initial data need only be imposed once. The Bianchi
identities and the conservation of energy and momentum tell us that, whether
or not the Einstein field equations are satisfied, we must have
\[
  \nabla_\nu\left(G^{\mu\nu}-8\pi G T^{\mu\nu}\right)=0.
\]
Let us apply this at \(x^0=t\). Imposing the initial-data constraints
\eqref{eq:7.5.1}, and determining the second derivatives from
Eq.~\eqref{eq:7.5.2}, the quantity in brackets will vanish everywhere at
\(x^0=t\), so this gives
\[
  \partial_0\left(G^{\mu0}-8\pi G T^{\mu0}\right)=0
  \qquad\text{at }x^0=t,
\]
and the fields computed at \(x^0=t+\delta t\) will therefore automatically
also satisfy the constraints \eqref{eq:7.5.1}. Thus this method of solving
the initial-value problem is one that can be programmed for an automatic
computer, once we find an initial metric at \(x^0=t\) that satisfies the
constraints \eqref{eq:7.5.1}.
